\section{Yamabe Problem Paper}

\setcounter{subsection}{1}
\subsection{Geometric and Analytic Preliminaries}
\label{sub_sec:geometric_analytic_preliminaries}

\begin{theorem}[Sobolev Embedding Theorems for $\R^n$]\leavevmode
  \label{thm:sobolev_embedding_rn}

  \begin{enumerate}
    \item Suppose
      \[%
        \frac{1}{r} = \frac{1}{q} - \frac{k}{n}
      .\]%
      Then, $L^q_k(\R^n) \hookrightarrow L^r(\R^n)$, continuously. In particular, for $q = 2$, $k = 1$, $r = p = 2n/(n - 2)$, we have the following Sobolev inequality:
      \begin{equation}\label{eq:sobolev_inequality}
        \|\phi\|^2_p \le \sigma_n \int_{\R^n} \left\lvert \nabla \phi\right\rvert^2 \dx
      ,\end{equation}
      for $\phi \in L^2_1(\R^n)$. We call $\sigma_n$ the \emph{best Sobolev constant}.

    \item Suppose $0 < \alpha < 1$, and
      \[%
        \frac{1}{q} \le \frac{k - \alpha}{n}
      .\]%
      Then, $L^q_k(\R^n) \hookrightarrow C^\alpha(\R^n)$, continuously.
  \end{enumerate}
\end{theorem}

\begin{computational-work}[\Cref{thm:sobolev_embedding_rn}]\leavevmode
  \label{cw:sobolev_embedding_rn}
\end{computational-work}

\begin{proof}[Proof of (i)]
\end{proof}

\begin{proof}[Proof of (ii)]
\end{proof}

\begin{theorem}[Sobolev Embedding Theorems for Compact Manifolds] \leavevmode
  \label{thm:sobolev_embedding_compact}

  \begin{enumerate}
    \item If
      \[%
        \frac{1}{r} \ge \frac{1}{q} - \frac{k}{n}
      ,\]%
      then $L^q_k(M) \hookrightarrow L^r(M)$, continuously.

    \item Suppose strict inequality holds in the above. Then, the inclusion $L^q_k(M) \subset L^r(M)$ is a compact operator.

    \item Suppose $0 < \alpha < 1$, and
      \[%
        \frac{1}{q} \le \frac{k - \alpha}{n}
      .\]%
      Then, $L^q_k(M) \hookrightarrow C^\alpha(M)$, continuously.
  \end{enumerate}
\end{theorem}

\begin{computational-work}[\Cref{thm:sobolev_embedding_compact}]\leavevmode
  \label{cw:sobolev_embedding_compact}
\end{computational-work}

\begin{proof}[Proof of (i)]
\end{proof}

\begin{proof}[Proof of (ii)]
\end{proof}

\begin{proof}[Proof of (iii)]
\end{proof}

\begin{theorem}[Aubin]
  \label{thm:aubin}

  Let $M$ be a compact Riemannian manifold with metric $g$, $p = 2n/(n - 2)$, and let $\sigma_n$ be the best Sobolev constant defined in~\cref{thm:sobolev_embedding_rn}. Then for every $\epsilon > 0$, there exists a constant $C_\epsilon$ such that for all $\phi \in C^\infty(M)$ such that
  \[%
    \|\phi\|^2_p \le (1 + \epsilon) \sigma_n \int_M |\nabla \phi|^2 \dd{V}_g + C_\epsilon \int_M \phi^2 \dd{V}_g
  .\]%
\end{theorem}

\begin{computational-work}[\Cref{thm:aubin}]\leavevmode
  \label{cw:aubin}

  \begin{enumerate}
    \item Showing $\|\phi\|_p^2 = \|\phi^2\|_{p/2}$:
      \[%
        \|\phi\|_p^2 = \left(\int_M |\phi|^p \dd{V}_g\right)^{2/p} = \left(\int_M |\phi^2|^{p/2} \dd{V}_g\right)^{1/(p/2)} = \|\phi^2\|_{p/2}
      .\]%

    \item Showing $\|\phi^2\|_{p/2} = \left\lVert\sum \psi_i^2 \phi^2\right\rVert_{p/2}$:

      Since $\sum_i \psi_i^2 = 1$, we have $|\phi^2| = |\sum_i \psi_i^2 \phi^2|$ pointwise, we get:
      \[%
        \|\phi^2\|_{p/2} = \left(\int_M |\phi|^p \dd{V}_g\right)^{1/(p/2)} = \left(\int_M \left|\sum_i \psi_i^2 \phi^2\right|^{p/2} \dd{V}_g\right)^{1/(p/2)} = \left\lVert\sum_i \psi_i^2 \phi^2\right\rVert_{p/2}
      .\]%

    \item Showing $\left\lVert \sum \psi_i^2 \phi^2\right\rVert \le \sum_i \left(\int_M |\psi_i \phi|^p \dd{V}_g\right)^{2/p}$:

      By the triangle inequality in $L^{p/2}$ and the identity $|\psi_i^2\phi^2|^{p/2} = |\psi_i\phi|^p$,
      \begin{align*}
        \left\lVert\sum_i \psi_i^2 \phi^2\right\rVert_{p/2} &= \left(\int_M \left|\sum_i \psi_i^2 \phi^2\right|^{p/2} \dd{V}_g\right)^{1/(p/2)} \\
                                                            &\le \sum_i \left(\int_M |\psi_i^2 \phi^2|^{p/2} \dd{V}_g\right)^{1/(p/2)} \\
                                                            &= \sum_i \left(\int_M |\psi_i \phi|^p \dd{V}_g\right)^{2/p}
      .\end{align*}

    \item Showing $\displaystyle\sum_i \left(\int_M |\psi_i\phi|^p\dd{V}_g\right)^{2/p} \le (1+\epsilon)^{n/2 \cdot 2/p} \sum_i \left(\int_{U_i} |\psi_i\phi|^p\dx\right)^{2/p}$:

      Since $\psi_i$ is supported on $U_i$, the integral over $M$ reduces to $U_i$. In normal coordinates on $U_i$, the eigenvalues of $g_{jk}$ lie in $\bigl[(1+\epsilon)^{-1},  1+\epsilon\bigr]$, so $\dd{V}_g \le (1+\epsilon)^{n/2}\dx$. Therefore,
      \[%
        \int_{U_i} |\psi_i\phi|^p \dd{V}_g \le (1+\epsilon)^{n/2} \int_{U_i} |\psi_i\phi|^p \dx
      ,\]%
      and raising to the power $2/p$ gives $(1+\epsilon)^{n/2 \cdot 2/p}$ per term. Since $p = 2n/(n-2)$, we have $n/2 \cdot 2/p = (n-2)/2 < 1$, so this exponent is harmless and may be absorbed into a relabelled $\epsilon$.

    \item Showing $|\nabla(\psi_i\phi)|^2 \le (1+\epsilon)\psi_i^2|\nabla\phi|^2 + (1+\epsilon^{-1})\phi^2|\nabla\psi_i|^2$:

      Expanding,
      \[%
        |\nabla(\psi_i\phi)|^2 = \psi_i^2|\nabla\phi|^2 + 2\psi_i\phi\langle\nabla\psi_i,\nabla\phi\rangle + \phi^2|\nabla\psi_i|^2
      .\]%
      Applying~\cref{thm:young_inequality} $2ab \le \epsilon a^2 + \epsilon^{-1}b^2$ with $a = \psi_i|\nabla\phi|$ and $b = \phi|\nabla\psi_i|$ to the cross term gives
      \begin{align*}
        |\nabla(\psi_i\phi)|^2 &\le \psi_i^2|\nabla\phi|^2 + \epsilon\psi_i^2|\nabla\phi|^2 + \epsilon^{-1}\phi^2|\nabla\psi_i|^2 + \phi^2|\nabla\psi_i|^2\\
                                &= (1+\epsilon)\psi_i^2|\nabla\phi|^2 + (1+\epsilon^{-1})\phi^2|\nabla\psi_i|^2
      .\end{align*}

    \item Relating the Euclidean and Riemannian gradients on $U_i$:

      In normal coordinates on $U_i$, the eigenvalues of $g_{jk}$ lie in $\bigl[(1+\epsilon)^{-1}, 1+\epsilon\bigr]$, so for any smooth function $f$,
      \[%
        |\nabla f|^2_0 \le (1+\epsilon) |\nabla f|^2_g, \qquad \dx \le (1+\epsilon)^{n/2}\dd{V}_g
      .\]%
      Therefore,
      \[%
        \int_{U_i} |\nabla(\psi_i\phi)|^2_0 \dx \le (1+\epsilon)^{1 + n/2} \int_{U_i} |\nabla(\psi_i\phi)|^2_g \dd{V}_g
      .\]%
  \end{enumerate}
\end{computational-work}

\begin{proof}
  Fix $\epsilon > 0$; we shall replace $(1+\epsilon)$ by a different name for the constant at the end, so we allow the factor to accumulate during the argument. Around each point of $M$ choose a coordinate neighbourhood on which the eigenvalues of $g_{jk}$ lie in $\bigl[(1+\epsilon)^{-1}, 1+\epsilon\bigr]$; this is possible by continuity of $g$ at every point. Extract a finite subcover $\{U_i\}$ and choose a subordinate partition of unity $\{\psi_i^2\}$ with $\sum_i \psi_i^2 = 1$ and each $\psi_i$ supported in $U_i$. By~\cref{cw:aubin}(i)--(iv),
  \begin{align*}
    \|\phi\|_p^2 = \|\phi^2\|_{p/2} = \Bigl\lVert\sum_i \psi_i^2\phi^2\Bigr\rVert_{p/2} &\le \sum_i \left(\int_M |\psi_i\phi|^p\dd{V}_g\right)^{2/p} \\
                                                                                        &\le (1+\epsilon)^{n/2 \cdot 2/p} \sum_i \left(\int_{U_i} |\psi_i\phi|^p\dx\right)^{2/p}
  .\end{align*}
  By~\cref{thm:sobolev_embedding_rn} applied to $\psi_i\phi \in C^\infty_c(\R^n)$,
  \[%
    \left(\int_{U_i} |\psi_i\phi|^p\dx\right)^{2/p} \le \sigma_n \int_{U_i} |\nabla(\psi_i\phi)|^2_0 \dx
  .\]%
  By~\cref{cw:aubin}(vi),
  \[%
    \int_{U_i} |\nabla(\psi_i\phi)|^2_0 \dx \le (1+\epsilon)^{1+n/2} \int_{U_i} |\nabla(\psi_i\phi)|^2_g \dd{V}_g
  .\]%
  By~\cref{cw:aubin}(v),
  \[%
    \int_{U_i} |\nabla(\psi_i\phi)|^2_g \dd{V}_g \le (1+\epsilon)\int_{U_i} \psi_i^2|\nabla\phi|^2_g \dd{V}_g + (1+\epsilon^{-1})\int_{U_i} \phi^2|\nabla\psi_i|^2_g \dd{V}_g
  .\]%
  Combining the previous steps and summing over $i$, we have:
  \begin{align*}
    \|\phi\|_p^2 &\le (1+\epsilon)^{n/2 \cdot 2/p} \cdot \sigma_n \cdot (1+\epsilon)^{1+n/2} \sum_i \left[(1+\epsilon)\int_{U_i}\psi_i^2|\nabla\phi|^2\dd{V}_g + (1+\epsilon^{-1})\int_{U_i}\phi^2|\nabla\psi_i|^2\dd{V}_g\right] \\
                  &= A(\epsilon) \sigma_n \int_M |\nabla\phi|^2\dd{V}_g + B(\epsilon)\int_M \phi^2\dd{V}_g
  ,\end{align*}
  where we used $\sum_i \psi_i^2 = 1$ to collapse the sums, and where
  \[%
    A(\epsilon) = (1+\epsilon)^{n/2 \cdot 2/p + 1 + n/2 + 1} \aand B(\epsilon) = A(\epsilon)\cdot(1+\epsilon^{-1})\sup_i\sup_{U_i}|\nabla\psi_i|^2
  .\]%
  Since $A(\epsilon) \to 1$ as $\epsilon \to 0$, for any prescribed $\epsilon' > 0$ we may choose $\epsilon$ small enough that $A(\epsilon) \le 1 + \epsilon'$. Setting $C_{\epsilon'} = B(\epsilon)$ completes the proof.
\end{proof}

\begin{theorem}[Local Elliptic Regularity]
  \label{thm:local_elliptic_regularity}

  Suppose $\Omega$ is an open set in $\R^n$, $\Delta$ is the Laplacian with respect to any metric on $\Omega$, and $u \in L^1_{\text{loc}}(\Omega)$ is a weak-solution to $\Delta u = f$.

  \begin{enumerate}
    \item If $f \in L^q_k(\Omega)$, then $u \in L^q_{k+2}(K)$ for any compact set $K \Subset \Omega$, and if $u \in L^q(\Omega)$, then
      \begin{equation}\label{eq:local_elliptic_regularity_thm}
        \|u\|_{L^q_{k+2}(K)} \le C\left(\|\Delta u\|_{L^q_k(\Omega)} + \|u\|_{L^q(\Omega})\right)
      \end{equation}

    \item If $f \in C^{k,\alpha}(\Omega)$, then $u \in C^{k+2,\alpha}(K)$ for any compact subset $K \Subset \Omega$, and if $u \in C^\alpha(\Omega)$, then
      \[%
        \|u\|_{C^{k+2,\alpha}(K)} \le C\left(\|\Delta u\|_{C^{k,\alpha}(\Omega)} + \|u\|_{C^\alpha(\Omega)}\right)
      .\]%
  \end{enumerate}
\end{theorem}

\begin{computational-work}[\Cref{thm:local_elliptic_regularity}]\leavevmode
  \label{cw:local_elliptic_regularity}
\end{computational-work}

\begin{proof}
  Expanding~\cref{eq:local_elliptic_regularity_thm}, we have:
  \begin{align*}
    \|u\|_{L^q_{k+2}(K)} &= \left(\sum_{i=0}^{k+2} \int_K \left\lvert\nabla^i u\right\rvert^q \dd{V}_g\right)^{1/q} \\
                          &\le C\left(\left(\sum_{i=0}^k \int_\Omega \left\lvert\nabla^i \left(\Delta u\right)\right\rvert^q \dd{V}_g\right)^{1/q} + \left(\int_\Omega |u|^q \dd{V}_g\right)^{1/q}\right)
  ,\end{align*}
  for some constant $C$ depending on $K$ and $\Omega$, and any metric $g$ on $\Omega$.
\end{proof}

\begin{proof}[Proof of (i)]
\end{proof}

\begin{proof}[Proof of (ii)]
\end{proof}

\begin{theorem}[Global Elliptic Regularity]
  \label{thm:global_elliptic_regularity}

  Let $M$ be a compact Riemannian manifold, and suppose $u \in L^1_{\text{loc}}(M)$ is a weak solution to $\nabla u = f$.

  \begin{enumerate}
    \item If $f \in L^q_k(M)$, then $u \in L^q_{k+2}(M)$, and
      \[%
        \|u\|_{q,k+2} \le C\left(\|\Delta u\|_{q,k} + \|\Delta u\|_q\right)
      .\]%
    \item If $f \in C^{k,\alpha}(M)$, then $u \in C^{k+2,\alpha}(M)$, and
      \[%
        \|u\|_{C^{k+2,\alpha}} \le C\left(\|\Delta u\|_{C^{k,\alpha}} + \|\Delta u\|_{C^\alpha}\right)
      .\]%
  \end{enumerate}
\end{theorem}

\begin{computational-work}[\Cref{thm:global_elliptic_regularity}]\leavevmode
  \label{cw:global_elliptic_regularity}
\end{computational-work}

\begin{proof}[Proof of (i)]
\end{proof}

\begin{proof}[Proof of (ii)]
\end{proof}

\begin{theorem}[Strong Maximum Principle]
  \label{thm:strong_maximum_principle}

  Suppose $h$ is a nonnegative smooth function on a connected manifold $M$, and $u \in C^2(M)$ satisfies $(\Delta + h)u \ge 0$. If $u$ attains its minimum $m \le 0$, then $u$ is constant on $M$.
\end{theorem}

\begin{computational-work}[\Cref{thm:strong_maximum_principle}]\leavevmode
  \label{cw:strong_maximum_principle}
\end{computational-work}

\begin{proof}
\end{proof}

\begin{proposition}[Weak Removable Singularities Theorem]
  \label{prop:weak_removable_singularities}

  Let $U$ be an open set in $M$ and $P \in U$. Suppose $u$ is a weak solution of $(\Delta + h)u = 0$ on $U - \{P\}$, with $h \in L^{n/2}(U)$ and $u \in L^q(U)$ for some $q > p/2 = n/(n - 2)$. Then $u$ satisfies $(\Delta + h)u = 0$ weakly on all of $U$.
\end{proposition}

\begin{computational-work}[\Cref{prop:weak_removable_singularities}]\leavevmode
  \label{cw:weak_removable_singularities}
\end{computational-work}

\begin{proof}
\end{proof}

\begin{theorem}[Existence of the Green Function]
  \label{thm:green_function}

  Suppose $M$ is a compact Riemannian manifold of dimension $n \ge 3$, and $h$ is a strictly positive smooth function on $M$. For each .
\end{theorem}

\begin{computational-work}[\Cref{thm:green_function}]\leavevmode
  \label{cw:green_function}
\end{computational-work}

\begin{proof}
\end{proof}

\subsection{The Model Case: The Sphere}

\subsection{Conformal Normal Coordinates}

\subsection{Stereographic Projections}

\subsection{The Test Function Estimate}

\subsection{General Relativity}

\subsection{Analysis on Asymptotically Flat Manifolds}

\subsection{The Positive Mass Theorem}

\subsection{Solution of the Yamabe Problem}
